\section{Nikodym and curved Kakeya consequences}\label{sec:consequences}

We apply the transfer theorems of Gao, Liu, and Xi to the maximal
estimate just proved. We state the local geometry explicitly so that
the domains and constant dependence remain visible.

\subsection{Nikodym maximal estimates on space forms}

Theorem~\ref{thm:main} also supplies the new input for the established
Kakeya-to-Nikodym transfer of Gao, Liu, and Xi
\citep[Section~2.2, Proposition~2.3]{GLX2025}. For a unit line segment
\(\gamma\subset\R^3\), write
\(T_\delta^{\mathrm{Euc}}(\gamma)=
\{y\in\R^3:\dist(y,\gamma)<\delta\}\), and define
\[
N_\delta^{\mathrm{Euc}}f(x)=
\sup_{\substack{\gamma\ni x\\ \gamma\text{ a unit line segment}}}
\frac1{\delta^2}\int_{T_\delta^{\mathrm{Euc}}(\gamma)}|f(y)|\,dy .
\]
This is the normalization in equation~(5) of \citet{GLX2025}.

For the manifold statement, let \((M,g)\) be a smooth three-dimensional
Riemannian manifold of constant sectional curvature. Fix open sets
\(U\Subset V\) with \(\overline V\) compact in a local chart that sends
geodesics to straight lines. Work at a fixed segment length within this
geometry, normalized to one by a fixed rescaling of \(g\). Choose
\(0<\delta_0<1\) small enough that the \(\delta_0\)-neighborhood of
every segment contained in \(U\) lies in \(V\). For \(0<\delta<\delta_0\),
put
\[
N_{\delta,U}^{g}f(x)=
\sup_{\substack{\gamma\ni x\\ \gamma\subset U}}
\frac1{\delta^2}\int_{T_\delta^{g}(\gamma)}|f(y)|\,dV_g(y),
\qquad x\in U,
\]
where the supremum is over unit geodesic segments, and
\(T_\delta^{g}(\gamma)=\{y\in M:d_g(y,\gamma)<\delta\}\).
An empty supremum is zero. This is a fixed local version of the
geodesic Nikodym operator in equation~(8) of \citet{GLX2025}.

\begin{corollary}[Nikodym maximal estimates in dimension three]
\label{cor:nikodym-space-forms}
For every \(1\le p\le3\) and \(\varepsilon>0\), set \(q=2p'\), with
\(q=\infty\) at \(p=1\). Then, for \(0<\delta<1\),
\[
\|N_\delta^{\mathrm{Euc}}f\|_{L^q(\R^3)}
\le C_{p,\varepsilon}\delta^{1-3/p-\varepsilon}
   \|f\|_{L^p(\R^3)}.
\]
For each fixed local geometry above and \(0<\delta<\delta_0\),
\[
\|N_{\delta,U}^{g}f\|_{L^q(U,dV_g)}
\le C_{p,\varepsilon,U,V,g}\delta^{1-3/p-\varepsilon}
   \|f\|_{L^p(V,dV_g)}.
\]
In particular, at \(p=q=3\) the operator norms are bounded by
\(C_\varepsilon\delta^{-\varepsilon}\) in Euclidean space and
\(C_{\varepsilon,U,V,g}\delta^{-\varepsilon}\) in the fixed local
geometry.
\end{corollary}

\begin{proof}
Write \(\widetilde K_\delta\) for the Kakeya operator of
\citet[Section~2.2]{GLX2025}, normalized by \(\delta^{-2}\) over
\(\delta\)-neighborhoods of unit segments. For \(0<\delta<1\), each
such neighborhood is contained in the union of three of the cylinders
\(T_\delta(a,\omega)\) used here. Consequently
\(\widetilde K_\delta f\le3\pi K_\delta f\), so
Theorem~\ref{thm:main} gives its \(L^3\)-to-\(L^3\) bound with every
positive power loss. Also
\[
\|\widetilde K_\delta f\|_{L^\infty(S^2)}
\le\delta^{-2}\|f\|_{L^1(\R^3)}.
\]
Interpolation for this sublinear operator, with
\(1/p=1-2\theta/3\) and \(1/q=\theta/3\), gives \(q=2p'\) and the
power \(-2(1-\theta)-\eta\theta=1-3/p-\eta\theta\) from an
\(L^3\) loss \(\delta^{-\eta}\). Since \(\eta>0\) is arbitrary,
including the two endpoint bounds gives \(K(3)\) in the notation of
\citet[equation~(6)]{GLX2025}.

Their Proposition~2.3 identifies \(K(3)\) with the Euclidean Nikodym
estimate \(N(3)\), and with the geodesic Nikodym estimate on a
constant-curvature manifold, in equations~(7) and~(9) respectively.
The negligible exponent loss in that transfer is absorbed by starting
with a smaller positive loss above. Its manifold comparison uses a
geodesic-straightening diffeomorphism whose bi-Lipschitz and volume
comparison constants are bounded on the fixed compact chart. Applying
that comparison to \(U\Subset V\) gives the stated local estimate,
with constants allowed to depend on this geometry.
\end{proof}

The constants in the manifold estimate depend on the fixed compact
straightening chart. The transfer uses the full maximal estimate of
Theorem~\ref{thm:main}, including its dependence on arbitrary shading
density.

\subsection{A local curved Kakeya consequence}

The same Euclidean maximal estimate also supplies the input for the
curved Kakeya transfer of Gao, Liu, and Xi for the following fixed
translation-invariant phase class. Fix a sufficiently small \(\rho>0\), put
\(\Omega=B_\rho^2\times(-\rho,\rho)\) and \(Y=B_\rho^2\), where
\(B_\rho^2\) is the open ball centered at the origin in \(\R^2\), and fix a
smooth real-valued phase on a neighborhood of
\(\overline{\Omega}\times\overline{Y}\) of the form
\[
\phi(x,t;y)=x\cdot y+\psi(t;y),\qquad \psi(0;y)=0.
\]
Assume the H\"ormander mixed-rank and curvature nondegeneracy conditions
\((\mathrm{H1})\)--\((\mathrm{H2})\) and Bourgain's condition there, in
the senses of \citet[Definitions~1.7, 1.10, and~1.11]{GLX2025}. For this phase, the mixed-rank condition is automatic, while the other
two conditions take the concrete form
\[
\det\nabla_y^2\partial_t\psi(t;y)\ne0,\qquad
\nabla_y^2\partial_t^2\psi(t;y)
=c(t;y)\nabla_y^2\partial_t\psi(t;y)
\]
for a scalar function \(c\). The phase and chart are fixed after any
shrinking needed for their local theorem.
For \(\omega,y\in B_\rho^2\) and \(0<\delta<\rho\), use the
gradient-defined tubes and normalization of their Definitions~1.8
and~4.5:
\begin{align*}
T^\phi_{\delta,y}(\omega)
&=\left\{(x,t)\in\Omega:
 \left|\nabla_y\phi(x,t;y)-\nabla_y\phi(\omega,0;y)\right|<\delta\right\},\\
\mathcal K^\phi_\delta f(y)
&=\sup_{\omega\in B_\rho^2}\frac1{\delta^2}
 \int_{T^\phi_{\delta,y}(\omega)}|f(x,t)|\,dx\,dt .
\end{align*}

\begin{corollary}[A local curved Kakeya maximal estimate]
\label{cor:curved-kakeya-local}
For the phase and chart fixed above, there is
\(0<\delta_{\mathrm{curv}}\le\rho\) such that, for every
\(1\le p\le3\) and \(\varepsilon>0\), there is a finite constant
\(C_{p,\varepsilon,\phi,\Omega,Y}\) for which
\[
\|\mathcal K^\phi_\delta f\|_{L^q(Y)}
\le C_{p,\varepsilon,\phi,\Omega,Y}
   \delta^{1-3/p-\varepsilon}\|f\|_{L^p(\Omega)}
\]
for every \(0<\delta<\delta_{\mathrm{curv}}\) and \(f\in L^p(\Omega)\),
where \(q=2p'\), with \(q=\infty\) at \(p=1\). In particular, at
\(p=q=3\),
\[
\|\mathcal K^\phi_\delta f\|_{L^3(Y)}
\le C_{\varepsilon,\phi,\Omega,Y}\delta^{-\varepsilon}
   \|f\|_{L^3(\Omega)}.
\]
\end{corollary}

\begin{proof}
The Euclidean comparison and interpolation in the preceding proof give
\(K(3)\) in the notation of \citet[equation~(6)]{GLX2025}. Their
Proposition~4.1 gives the equivalent assertion \(\mathcal K(3)\) for
the phase class above, which is exactly the displayed \(p,q\) range.
Their straightening theorem is local
\citep[Theorem~1.12]{GLX2025}; on the fixed chart its spatial and
direction-parameter comparison constants may depend on the phase and
the chart. This gives the stated estimates.
\end{proof}

The constants are allowed to depend on the phase and localization,
which are fixed independently of \(\delta\).
